\section{Couplings and CP violation as geometric normalization}
\label{sec:couplings_cp}

This section records two interface points: (i) coupling constants as geometric normalization costs, and (ii) CP violation as a CP-odd phase-space volume with discrete multiplicity.
We treat the closed expressions as interface-level normalizations: they follow from explicitly stated axioms and canonical geometric data, and their mismatch to scheme-/scale-dependent experimental conventions is recorded as a matching-layer factor.

\paragraph{Detailed constructions (supplement).}
The explicit geometric constructions, bounded-complexity enumerations, and CKM reconstruction tables are recorded in Appendix~\ref{app:closure_audit_details}.

\subsection{Summary table: closed values vs.\ CODATA/PDG}
\label{subsec:summary_table}

\begin{table}[t]
\centering
\small
\setlength{\tabcolsep}{6pt}
\renewcommand{\arraystretch}{1.15}
\begin{tabular}{lccc}
\toprule
Quantity & closed value & reference (CODATA/PDG) & $\log(\mathrm{closed}/\mathrm{ref})$ \\
\midrule
$\alpha_{\mathrm{em}}^{-1}$ (low energy) & $4\pi^3+\pi^2+\pi$ & $137.035999084$ (CODATA 2022) & $+2.22\times 10^{-6}$ \\
$\alpha^{-1}(\mu_Z)$ & $13\pi^2$ & $127.955$ (PDG) & $+2.73\times 10^{-3}$ \\
$\sin^2\theta_W(\mu_Z)$ & $3/13$ & $0.23122$ (PDG, $\overline{\mathrm{MS}}$) & $-1.95\times 10^{-3}$ \\
$J$ (CKM) & $1/(11\pi^7)$ & $3.00\times 10^{-5}$ (PDG) & $+3.31\times 10^{-3}$ \\
\bottomrule%
\end{tabular}
\caption{Quantitative rigidity targets of the closed model in this paper, together with standard reference values and signed log mismatches in the audit norm. For small deviations, $\log(\mathrm{closed}/\mathrm{ref})\approx (\mathrm{closed}-\mathrm{ref})/\mathrm{ref}$. Domain sizes, uniqueness gaps, uncertainty-robustness checks, and counterfactual baselines for the associated bounded-complexity closures are recorded in Appendix~\ref{app:generated_tables} (Tables~\ref{tab:audit_closure_metrics}--\ref{tab:audit_counterfactual}) \cite{PDGRPP2024,CODATA2022RMP}. Rows are generated by \texttt{scripts/exp\_quant\_summary.py}.}
\label{tab:quant_summary}
\end{table}

\paragraph{Interpretation of mismatch.}
The closed expressions in Table~\ref{tab:quant_summary} are not free fit parameters: they are rigidity targets selected by explicit bounded-complexity closure rules.
The mismatch $\log(\mathrm{closed}/\mathrm{ref})$ is recorded as a protocol-level matching factor between an idealized closed normalization and a scheme-/scale-dependent reference convention.
This plays the same structural role as the matching-layer depth shift $\Delta r$ used later for masses (Appendix~\ref{app:mass_rigidity_details}).
The corresponding finite rigidity enumerations (top candidates and gaps within the declared domains) are recorded in Appendix~\ref{app:generated_tables} (Tables~\ref{tab:alpha_coeff_rigidity}--\ref{tab:jarlskog_pi_rigidity}).

\begin{table}[t]
\centering
\scriptsize
\setlength{\tabcolsep}{6pt}
\renewcommand{\arraystretch}{1.15}
\resizebox{\textwidth}{!}{%
\begin{tabular}{lrrrrr}
\toprule
Quantity & closed value & reference & $\sigma$ (audit) & error & $|\mathrm{error}|/\sigma$ \\
\midrule
$\alpha_{\mathrm{em}}^{-1}$ & $137.0363$ & $137.036$ & $2.1\times 10^{-8}$ & $3.047\times 10^{-4}$ & $1.451\times 10^{4}$ \\
$\alpha^{-1}(\mu_Z)$ & $128.30486$ & $127.955$ & $0.01$ & $0.34985721$ & $34.985721$ \\
$\sin^2\theta_W(\mu_Z)$ & $0.23076923$ & $0.23122$ & $3\times 10^{-5}$ & $-4.508\times 10^{-4}$ & $15.025641$ \\
$J$ (CKM) & $3.01\times 10^{-5}$ & $3\times 10^{-5}$ & $1.5\times 10^{-6}$ & $9.943\times 10^{-8}$ & $0.066283647$ \\
$|V_{us}|$ & $0.2236068$ & $0.2243$ & $5\times 10^{-4}$ & $-6.932\times 10^{-4}$ & $1.3864045$ \\
$|V_{cb}|$ & $0.043810715$ & $0.0422$ & $8\times 10^{-4}$ & $0.0016107147$ & $2.0133934$ \\
$|V_{ub}|$ & $0.0039504097$ & $0.00394$ & $3.6\times 10^{-4}$ & $1.041\times 10^{-5}$ & $0.028915787$ \\
$\sin^2\theta_{12}$ (PMNS) & $0.30769231$ & $0.307$ & $0.013$ & $6.923\times 10^{-4}$ & $0.053254438$ \\
$\sin^2\theta_{23}$ (PMNS) & $0.54545455$ & $0.545$ & $0.021$ & $4.545\times 10^{-4}$ & $0.021645022$ \\
$\sin^2\theta_{13}$ (PMNS) & $0.021286236$ & $0.0218$ & $7\times 10^{-4}$ & $-5.138\times 10^{-4}$ & $0.73394821$ \\
$\delta$ (PMNS) [deg] & $195$ & $195$ & $30$ & $0$ & $0$ \\
\bottomrule%
\end{tabular}
}%
\caption{Sigma-normalized mismatch summary for selected key targets. Here $\sigma$ denotes the explicit uncertainty scales used by the audit scripts (either quoted uncertainties or deterministic stress-test scales; see \texttt{scripts/common\_constants.py}). This table is provided for interpretive convenience and does not replace the matching-layer interpretation when reference conventions are scheme/scale dependent or when the quoted mismatch is much larger than the reference uncertainty. Rows are generated by \texttt{scripts/exp\_sigma\_summary.py}.}
\label{tab:sigma_summary}
\end{table}


